\Needspace{8\baselineskip}
\Article{280}
Suppose that the binary form
\[
att+2btu+cuu\ldots\varphi
\]
\SourceLeaf{324}
whose determinant $=D$, is represented by the ternary form $f$, with
indeterminates $x,x',x''$, by setting
\[
x=mt+nu,\qquad x'=m't+n'u,\qquad x''=m''t+n''u
\]
and that $f$ has the adjoint form $F$, with indeterminates $X,X',X''$.
Then a calculation readily confirms (if we use separate letters for the
coefficients of the forms $f,F$), or Article 268.\ II.\ directly shows,
that the number $D$ is represented by $F$ on setting
\[
X=m'n''-m''n',\qquad X'=m''n-mn'',\qquad X''=mn'-m'n
\]
This representation of the number $D$ can conveniently be called
\textit{adjoint} to the representation of the form $\varphi$ by $f$.
If the values of $X,X',X''$ have no common divisor, we will, for brevity,
call this representation of $D$ \textit{proper}; otherwise,
\textit{improper}. We will give the same names to the representation of
the form $\varphi$ by $f$ to which that representation of $D$ is
adjoint. Finding all proper representations of the number $D$ by
the form $F$ now rests on the following points:

I.\ Every representation of $D$ by $F$ can be derived from some
representation of some form of determinant $D$ by the form $f$,
\textit{i.e.}, is adjoint to such a representation.

For let any representation of $D$ by $F$ be given by $X=L$, $X'=L'$,
$X''=L''$. By the lemma in the preceding article, choose $m,m',m'',n,n',n''$ so that
\[
m'n''-m''n'=L,\qquad m''n-mn''=L',\qquad mn'-m'n=L''
\]
and let $f$ be transformed by the substitution
\[
x=mt+nu,\qquad x'=m't+n'u,\qquad x''=m''t+n''u
\]
into the binary form $\varphi=att+2btu+cuu$. It is then easy to see that $D$ will be
the determinant of the form $\varphi$, and that its representation by $f$
will have the proposed representation of $D$ by $F$ as its adjoint.

\textit{Example.} Let $f=xx+x'x'+x''x''$, so that $F=-XX-X'X'-X''X''$;
let $D=-209$, with the representation by $F$ given by $X=1$, $X'=8$, $X''=12$.
\SourceLeaf{325}
This gives $m,m',m'',n,n',n''$ the values $-20,1,1,-12,0,1$, respectively, and gives $\varphi=402tt+482tu+145uu$.

II.\ If $\varphi,\chi$ are properly equivalent binary forms, every representation
of $D$ by $F$ that is adjoint to some representation of the form
$\varphi$ by $f$ will also be adjoint to some representation of the form
$\chi$ by $f$.

Let $p,q$ be the indeterminates of the form $\chi$; let $\varphi$ be
transformed into $\chi$ by the proper substitution $t=\alpha p+\beta q$, $u=\gamma p+\delta q$,
and let some representation of the form $\varphi$ by $f$ be given by
\[
x=mt+nu,\qquad x'=m't+n'u,\qquad x''=m''t+n''u\ldots(R)
\]
It is then easy to see that, if we set
\[
\begin{aligned}
\alpha m+\gamma n&=g,&\quad \alpha m'+\gamma n'&=g',&\quad \alpha m''+\gamma n''&=g''\\
\beta m+\delta n&=h,&\quad \beta m'+\delta n'&=h',&\quad \beta m''+\delta n''&=h''
\end{aligned}
\]
the form $\chi$ will be represented by $f$ on setting
\[
x=gp+hq,\qquad x'=g'p+h'q,\qquad x''=g''p+h''q\ldots(R')
\]
and a calculation gives (since $\alpha\delta-\beta\gamma=1$)
\[
g'h''-g''h'=m'n''-m''n',\quad
g''h-gh''=m''n-mn'',\quad
gh'-g'h=mn'-m'n
\]
\textit{i.e.}, the representations $R,R'$ have the same representation
of $D$ by $F$ as their adjoint.

Thus, in the preceding example, the form $\varphi$ is found to be equivalent to
$\chi=13pp-10pq+18qq$, into which it is transformed by the proper substitution
$t=-3p+q$, $u=5p-2q$. This gives the representation of the form $\chi$ by $f$
as $x=4q$, $x'=-3p+q$, $x''=2p-q$, from which we obtain the same representation
of the number $-209$ that we started with.

III.\ Finally, suppose that two binary forms $\varphi,\chi$ of determinant $D$,
with indeterminates $t,u$; $p,q$, can be represented by $f$,
and that the same proper representation of $D$ by $F$ is adjoint
to some representation of one and to some representation of the other.
Then these forms must be properly equivalent. Suppose that $\varphi$ is
represented by $f$ on setting
\[
x=mt+nu,\qquad x'=m't+n'u,\qquad x''=m''t+n''u
\]
\SourceLeaf{326}
and $\chi$ on setting
\[
x=gp+hq.\qquad x'=g'p+h'q,\qquad x''=g''p+h''q
\]
and that
\[
\begin{aligned}
m'n''-m''n'&=g'h''-g''h'=L\\
m''n-mn''&=g''h-gh''=L'\\
mn'-m'n&=gh'-g'h=L''
\end{aligned}
\]
Choose integers $l,l',l''$ so that $Ll+L'l'+L''l''=1$, and set
\[
\begin{aligned}
n'l''-n''l'&=M,&\quad n''l-nl''&=M',&\quad nl'-n'l&=M''\\
l'm''-l''m'&=N,&\quad l''m-lm''&=N',&\quad lm'-l'm&=N''
\end{aligned}
\]
Finally, set
\[
\begin{aligned}
gM+g'M'+g''M''&=\alpha,&\quad hM+h'M'+h''M''&=\beta\\
gN+g'N'+g''N''&=\gamma,&\quad hN+h'N'+h''N''&=\delta
\end{aligned}
\]
We readily obtain
\[
\begin{aligned}
\alpha m+\gamma n&=g-l(gL+g'L'+g''L'')=g\\
\beta m+\delta n&=h-l(hL+h'L'+h''L'')=h
\end{aligned}
\]
and, in the same way,
\[
\alpha m'+\gamma n'=g',\quad \beta m'+\delta n'=h',\quad
\alpha m''+\gamma n''=g'',\quad \beta m''+\delta n''=h''
\]
This shows. $mt+nu$, $m't+n'u$, $m''t+n''u$ are taken by the substitution
\[
t=\alpha p+\beta q,\quad u=\gamma p+\delta q\ldots(S)
\]
into $gp+hq$, $g'p+h'q$, $g''p+h''q$, respectively. So $\varphi$ is clearly
taken by the substitution $S$ into the same form as $f$ is on setting
\[
x=gp+hq,\qquad x'=g'p+h'q,\qquad x''=g''p+h''q
\]
and therefore into the form $\chi$, to which it is thus equivalent.
Finally, making the required substitutions readily gives
\[
\alpha\delta-\beta\gamma=(Ll+L'l'+L''l'')^2=1
\]
so the substitution $S$ is proper, and the forms $\varphi,\chi$ are properly equivalent.
\SourceLeaf{327}

These observations give the following rules for finding all proper
representations of $D$ by $F$: Work out all classes of binary
forms of determinant $D$, and choose one form at will from each.
Find all proper representations of each of these forms by $f$
(discarding those that cannot be represented by $f$), and from each
of these representations derive the representations of the number $D$
by $F$. From I and II it is clear that this gives all possible proper
representations, so the solution is complete; from III, that transformations
of forms from different classes certainly produce different representations.

\Needspace{8\baselineskip}
\Article{281}
Finding the \textit{improper} representations of a given number $D$
by the form $F$ is easily reduced to the preceding case. For if $D$
is divisible by no square (other than $1$), clearly there are no such
representations at all. Otherwise, if $D$ is divisible by the squares $\lambda\lambda,\mu\mu,\nu\nu$
and so on, all improper representations of $D$ by $F$ are found
by working out all proper representations of the numbers $\frac{D}{\lambda\lambda}$,
$\frac{D}{\mu\mu}$, $\frac{D}{\nu\nu}$, and so on by the same form, and multiplying the values
of the indeterminates by $\lambda,\mu,\nu$ and so on, respectively.

In this way, finding all representations of a given number by a given
ternary form \textit{that is adjoint to some ternary form} depends on the
second problem. The remaining cases reduce to this case, which at first
may seem to have rather limited scope, as follows. Let $D$ be the
number to be represented by the form
$\left(\begin{smallmatrix}g,g',g''\\h,h',h''\end{smallmatrix}\right)$,
with determinant $\Delta$, whose adjoint form is
$\left(\begin{smallmatrix}G,G',G''\\H,H',H''\end{smallmatrix}\right)=f$.
The form adjoint to this in turn will be
$\left(\begin{smallmatrix}\Delta g,\Delta g',\Delta g''\\\Delta h,\Delta h',\Delta h''\end{smallmatrix}\right)=F$,
and clearly the representations of the number $\Delta D$ by $F$
(whose determination depends on what precedes) are exactly identical
with the representations of the number $D$ by the proposed form.
--- Further, when all coefficients of the form $f$ have a common
divisor $\mu$, all coefficients of the form $F$ are clearly
divisible by $\mu\mu$. So $\Delta D$ must also be divisible by $\mu\mu$
(otherwise there would be no representations); and the representations
of the number $D$ by the proposed form coincide with the representations
of the number $\frac{\Delta D}{\mu\mu}$ by the form obtained from $F$ by dividing each
coefficient by $\mu\mu$. This form is adjoint to the one obtained from
$f$ by dividing each coefficient by $\mu$.

Finally, we note that this solution of the first problem does not apply
in the single case where
\SourceLeaf{328}
$D=0$; for here not all binary forms of determinant $D$ fall into
a finite number of classes. We will solve this case below on other principles.

